\subsection{Plan de Aceptación del Producto - Comprobar el producto con quienes lo usarán [Pendientes: 1. Hablar con la directora de Bomberos de Bogotá para acordar su participación en la aceptación del producto. 2. Acordar criterios medibles, métodos de comprobación, tolerancias y responsables.]}\label{sec:spmp-8-3}
La aceptación se realiza en dos niveles. En cada iteración, el \textit{Product Owner} acepta los incrementos que cumplen la definición de terminado de la tabla~\ref{tab:aceptacion} y los presenta en la revisión con el director. Al cierre de cada fase, y en la entrega final, el director del proyecto acepta formalmente los entregables E1 a E7 (sección~6.4) con los criterios de la misma tabla. Los representantes de la UAECOB no aprueban entregables: aportan la evaluación externa PC07, cuyos resultados forman parte de los criterios de E5.
 
Cada entregable presentado recibe una de tres decisiones: aceptado; aceptado con observaciones, que se corrigen en la iteración siguiente, o no aceptado, en cuyo caso se corrige y se presenta de nuevo. La decisión, sus mediciones y la versión evaluada se registran en E6 y en el tablero del proyecto. Un defecto es crítico si impide completar un caso de uso o presenta resultados en una ubicación geográfica incorrecta. Todo cambio a estos criterios sigue el procedimiento de la sección~6.6.
 
\begin{longtable}{|p{0.08\textwidth}|p{0.45\textwidth}|p{0.25\textwidth}|p{0.09\textwidth}|}
\caption{Criterios de aceptación de los entregables. Fuente: elaboración propia.}\label{tab:aceptacion}\\
\hline
\textbf{ID} & \textbf{Criterios de aceptación} & \textbf{Técnica y herramienta} & \textbf{Sección}\tabularnewline\hline
\endfirsthead
\hline
\textbf{ID} & \textbf{Criterios de aceptación} & \textbf{Técnica y herramienta} & \textbf{Sección}\tabularnewline\hline
\endhead
E1 & \raggedright El 100~\% de los requisitos comprometidos tiene identificador, criterio de aceptación y relación con al menos un caso de prueba en la matriz de trazabilidad. La revisión no deja hallazgos críticos abiertos. & \raggedright Revisión con lista de comprobación (PC01); matriz de trazabilidad. & \raggedright 10.1 y 12.5\tabularnewline\hline
E2 & \raggedright Cada caso de uso comprometido tiene su diagrama de secuencia, y cada componente del diagrama de componentes corresponde a un módulo del código. La revisión no deja hallazgos críticos abiertos. & \raggedright Revisión de modelos (PC01) y auditoría de artefactos (PC06); PlantUML. & \raggedright 12.3 y 12.5\tabularnewline\hline
E3 & \raggedright Al repetir el caso piloto con la misma configuración, versión de FDS y número de procesos, el área quemada y la posición del frente en los tiempos de control difieren menos de 1~\%. El informe del piloto reporta el costo $c_1$, la memoria y el tamaño de salida medidos. & \raggedright Reejecución de comprobación (PC02); \texttt{fdsreader}, \texttt{sacct} y \texttt{du}. & \raggedright 9.1 y 12.5\tabularnewline\hline
E4 & \raggedright La biblioteca contiene el número de escenarios definido tras el piloto (sección~9.1). El 100~\% de los escenarios tiene metadatos completos (ignición, meteorología, configuración y versión de FDS) y una suma de verificación SHA-256 comprobada, y el catálogo se carga sin errores en la aplicación. & \raggedright Verificación de datos (PC02) y auditoría (PC06); \textit{script} de verificación en Python. & \raggedright 9.1, 12.3 y 12.5\tabularnewline\hline
\raggedright E5 por iteración & \raggedright Definición de terminado de cada incremento: cumple los criterios de aceptación de su caso de uso en E1, el 100~\% de sus pruebas unitarias se aprueba, su solicitud de integración fue revisada por un integrante distinto del autor y no tiene defectos críticos abiertos. & \raggedright Pruebas unitarias (PC03) con pytest; revisión en GitHub; demostración en la revisión con el director. & \raggedright 8.1, 12.3 y 12.5\tabularnewline\hline
\raggedright E5 final & \raggedright El 100~\% de los casos de prueba de sistema de CU01 a CU06 se aprueba, incluidos los casos controlados de agregación. El ejecutable funciona en Windows 10 y 11 sin instalaciones previas. Ninguna de diez consultas de prueba supera 5 minutos, desde el ingreso del perímetro y la meteorología hasta la animación en Smokeview. Con los evaluadores de PC07, el puntaje SUS promedio es de al menos 68 y al menos el 80~\% de las tareas se completa sin ayuda. & \raggedright Pruebas de integración y sistema (PC04 y PC05); cronometraje; cuestionario SUS~\cite{brooke1996sus,lewis2018sus} y observación de tareas (PC07) con Google Forms. & \raggedright 8.2, 9.1 y 12.5\tabularnewline\hline
E6 & \raggedright El 100~\% de los casos de prueba ejecutados tiene resultado registrado y vinculado a la versión evaluada, y todos los hallazgos de la auditoría están cerrados o justificados. & \raggedright Auditoría interna (PC06); registros de pruebas. & \raggedright 10.2 y 12.5\tabularnewline\hline
E7 & \raggedright Un integrante que no escribió el manual instala la aplicación y completa una consulta siguiéndolo, sin ayuda. Otro integrante reproduce la construcción del ejecutable con la guía técnica. El director aprueba el informe final tras la demostración de CU01 a CU06. & \raggedright Prueba de instalación en un equipo limpio; revisión del director. & \raggedright 10.3 y 11\tabularnewline\hline
\end{longtable}
 
La evaluación PC07 requiere al menos tres evaluadores de la UAECOB, coordinados por el \textit{Product Owner}, que realizan tareas representativas de CU01 a CU06 y luego responden el cuestionario SUS. El 68 corresponde al puntaje promedio de referencia de la escala~\cite{lewis2018sus}. Si la UAECOB no aporta tres evaluadores, PC07 se reporta como no realizado y la aceptación de E5 se decide con los demás criterios.
 